% A l'original, el fragment «com més petita… del microscopi» és subratllat (paraula a paraula perquè pugui partir la línia).
Per a fer observacions, els microscopis òptics fan servir fotons i els microscopis electrònics utilitzen electrons. El poder de resolució d'un microscopi és la capacitat que té per a diferenciar com a punts separats dos punts molt propers i està determinat (en part) per la longitud d'ona de la radiació emprada, de tal manera que \underline{com} \underline{més} \underline{petita} \underline{és} \underline{la} \underline{longitud} \underline{d'ona} \underline{de} \underline{la} \underline{radiació,} \underline{més} \underline{gran} \underline{és} \underline{la} \underline{resolució} \underline{del} \underline{microscopi}.

\figura[0.35]{fig1.png}

\begin{apartat}{1}
Calculeu l'energia dels fotons utilitzats en un microscopi òptic de llum visible de 400~nm de longitud d'ona. Quina és la quantitat de moviment d'aquests fotons?
\end{apartat}

\begin{apartat}{1}
Fem servir un microscopi electrònic en què els electrons que ens permeten fer l'observació són accelerats per una diferència de potencial, de manera que assoleixen una quantitat de moviment de $3,31\times 10^{-25}\ \mathrm{kg\,m\,s^{-1}}$. Calculeu la relació que hi ha entre el poder de resolució d'aquest microscopi electrònic i el del microscopi òptic de l'apartat anterior. Quin dels dos microscopis té més poder de resolució?
\end{apartat}

\dades{%
  Velocitat de la llum, $c = 3,00\times 10^{8}\ \mathrm{m\,s^{-1}}$.\\
  Constant de Planck, $h = 6,63\times 10^{-34}\ \mathrm{J\,s}$.}
